\section{Conclusion} \label{sec:conclusion}
We investigated whether Randers/Finsler-style asymmetric retrieval scores improve domain-generalizable person ReID on top of BAU-style alignment-uniformity training. Our primary finding is that the \textbf{Arm~1b} recipe (eq.~\ref{eq:total_loss}; Table~\ref{tab:aux_loss_sweep}), with drift-subspace variance regularization rather than camera-mined domain triplets or domain-prototype routing, yields stable asymmetric retrieval on M+MS+CS$\rightarrow$CUHK03.  The domain-prototype head, while principled, does not improve over the per-sample residual MLP when $\mathcal{L}_{\mathrm{dcc}}$ already bounds cross-camera drift dispersion. 

\paragraph{Scope and limitations.}
Multi-target evaluation across all three Protocol~2 and 3 splits is reported in Table~\ref{tab:template_multiprotocol_bau}; multi-seed analysis (Table~\ref{tab:multiseed}, five seeds) yields mAP $= 43.8 \pm 0.3\,\mathrm{pp}$ on M+MS+CS$\rightarrow$C3, confirming stable training. Finsler BAU outperforms the matched Euclidean baseline by $+1.9\,\mathrm{pp}$ mAP on M+MS+CS$\rightarrow$C3 (Protocol~2) but is within $0.3\,\mathrm{pp}$ on all other splits and under Protocol~3, indicating that the benefit of asymmetric geometry is split- and augmentation-dependent. The controlled toy study (Sec.~\ref{subsec:controlled_corruption}) validates the asymmetric scoring mechanism independently of the full training stack: $\mathcal{L}_{\mathrm{mono}}$ reliably encodes corruption severity in a learned scalar projection; the Randers correction produces the predicted directional gap; and the balanced bidirectional evaluation confirms positive directional asymmetry at matched counts (statistics in Table~\ref{tab:toy_balanced_baseline} and~\ref{tab:toy_balanced_trained}, disabled-loss control in Appendix~\ref{sec:suppl_m2b_ablation}). Toy mAP deltas from Randers alone remain small ($N\!=\!50$, single seed); the bidirectional triplet gain ($+1.0\,\mathrm{pp}$ mAP) in the main system comes through training geometry.


\paragraph{Future work.}
Several directions remain open. \textbf{Cross-modal retrieval} (RGB-IR, RGB-depth) is the most natural next application: sensor-induced asymmetry is large, fixed in direction, and the domain prior is exactly known: conditions where the domain-prototype head and a larger drift coherence weight may be appropriate. \textbf{ViT/Swin backbones} may expose qualitatively different drift behavior: patch-level representations do not pool global statistics the same way as GeM-pooled ResNet50 features, and the backbone coupling identified in Sec.~\ref{subsec:structured_embedding} will take a different form in the attention regime. \textbf{Drift-only triplet mining} is unexplored: the drift subspace currently receives no direct identity-discriminative or domain-discriminative triplet signal of its own, relying entirely on the bidirectional unified triplet and $\mathcal{L}_{\mathrm{dcc}}$. Mining within the drift subspace alone may expose whether drift carries retrievable structure beyond variance regularization. \textbf{Capacity attribution} is unresolved: Finsler BAU uses a 4096-d embedding vs.\ the 2048-d Euclidean baseline; whether the mAP gain is attributable to asymmetric geometry or to additional model capacity remains to be tested with a matched 4096-d Euclidean baseline. \textbf{Statistical disentanglement} requires breaking the shared-backbone constraint: a dedicated identity encoder and a separate drift encoder (e.g., DSN-style~\cite{bousmalis2016domain}) would sever the Jacobian coupling identified in Sec.~\ref{subsec:structured_embedding}. We also flag using the Randers distance itself as a direct supervision signal for drift scaling (rather than via Gram-Schmidt plus barrier), citing \cite{bao2000introduction,randers1941geometry}, as a direction not pursued here.
